%_______________
\eocesolch{Distribusi variabel acak}



%_______________
\begin{multicols}{2}

% 1

\eocesol{(a)~8.85\%.
(b)~6.94\%.
(c)~58.86\%.
(d)~4.56\%. \\
\FigureFullPath[Distribusi normal yang berpusat di 0 ditampilkan dengan ekor kiri kecil pada dan di bawah lokasi berlabel -1.35 diberi arsiran.]{0.23}{ch_distributions/figures/eoce/area_under_curve_1/zltNeg}
\FigureFullPath[Distribusi normal yang berpusat di 0 ditampilkan dengan ekor kanan kecil pada dan di atas lokasi berlabel 1.48 diberi arsiran.]{0.23}{ch_distributions/figures/eoce/area_under_curve_1/zgtPos}
\FigureFullPath[Distribusi normal yang berpusat di 0 ditampilkan dengan daerah tengah diberi arsiran. Ekor kiri yang besar sampai sedikit di bawah rata-rata dan ekor kanan yang kecil tetap tidak diarsir.]{0.23}{ch_distributions/figures/eoce/area_under_curve_1/zBet}
\FigureFullPath[Distribusi normal yang berpusat di nol ditampilkan dengan kedua ekor di bawah -2 dan di atas 2 diberi arsiran.]{0.23}{ch_distributions/figures/eoce/area_under_curve_1/zgtAbs}}

% 3

\eocesol{(a)~Verbal: $N(\mu = 151, \sigma = 7)$, Kuantitatif: $N(\mu = 153, \sigma = 7.67)$.
(b)~$Z_{VR} = 1.29$, $Z_{QR} = 0.52$. \\
\FigureFullPath[Ditampilkan sebuah distribusi normal dengan 2 garis vertikal yang diberi tanda khusus. Garis pertama terletak sedikit di atas rata-rata distribusi pada Z sama dengan 0.52 dan diberi label “QR” untuk Penalaran Kuantitatif. Garis kedua terletak lebih jauh di atas rata-rata pada Z sama dengan 1.29 dan diberi label “VR” untuk Penalaran Verbal.]{0.3}{ch_distributions/figures/eoce/GRE_intro/gre_intro.pdf} \\
(c)~Skornya berada 1.29 simpangan baku di atas rata-rata pada bagian
Penalaran Verbal dan 0.52 simpangan baku di atas rata-rata pada bagian
Penalaran Kuantitatif.
(d)~Ia berprestasi lebih baik pada bagian Penalaran Verbal karena skor-Z-nya
lebih tinggi pada bagian tersebut.
(e)~$Perc_{VR} = 0.9007 \approx 90\%$, $Perc_{QR} = 0.6990 \approx 70\%$.
(f)~$100\% - 90\% = 10\%$ memperoleh skor lebih tinggi daripadanya pada VR,
dan $100\% - 70\% = 30\%$ memperoleh skor lebih tinggi daripadanya pada QR.
(g)~Kita tidak dapat membandingkan skor mentah karena kedua bagian menggunakan
skala yang berbeda. Perbandingan skor persentil lebih tepat ketika membandingkan
kinerjanya dengan peserta lain.
(h)~Jawaban bagian (b)--(c) tidak berubah karena skor-Z dapat dihitung untuk
distribusi yang tidak normal. Namun, kita tidak dapat menjawab bagian~(d)--(f)
karena tanpa model normal kita tidak dapat menggunakan tabel peluang normal
untuk menghitung peluang dan persentil.}

% 5

\eocesol{(a)~$Z = 0.84$, yang bersesuaian dengan skor sekitar 159 pada QR.
(b)~$Z = -0.52$, yang bersesuaian dengan skor sekitar 147 pada VR.}

% 7

\eocesol{(a)~$Z = 1.2$, $P(Z > 1.2) = 0.1151$. \\
(b)~$Z= -1.28 \to 70.6\degree$F atau lebih rendah.}

% 9

\eocesol{(a)~$N(25, 2.78)$.
(b)~$Z = 1.08$, $P(Z > 1.08) = 0.1401$.
(c)~Jawabannya sangat berdekatan karena yang berubah hanyalah satuan.
(Jawabannya sedikit berbeda karena 28\degree C sama dengan 82.4\degree F,
bukan tepat 83\degree F.)
(d)~Karena $IQR = Q3 - Q1$, pertama-tama kita perlu mencari $Q3$ dan $Q1$,
kemudian menghitung selisihnya. Ingat bahwa $Q3$ adalah persentil ke-$75$
dan $Q1$ adalah persentil ke-$25$ suatu distribusi.
Dengan kuantil yang belum dibulatkan, Q1 = 23.1264\degree C,
Q3 = 26.8736\degree C, sehingga IQR = 3.7472\degree C,
atau sekitar 3.75\degree C.}
